% =============================================================================
% UNIFIED TRAINING OBJECTIVE - CITA Paper in ALKALI Format
% =============================================================================

\vspace{-2mm}
\section{Unified Training Objective}
\label{sec:unified_loss}

CITA combines three complementary loss components to achieve instruction-conditioned alignment, drawing on insights from supervised fine-tuning~\cite{wei2022finetuned,sanh2022multitask}, contrastive preference learning~\cite{rafailov2023direct,chen2020simple}, and KL-regularized optimization~\cite{schulman2017proximal,korbak2022rl}. The key insight is that each component serves a distinct purpose, and their combination---particularly the \textit{mandatory} KL term---enables stable behavioral switching.

\vspace{-2mm}
\subsection{Loss Formulation}
\label{sec:loss_formulation}

The unified CITA objective is:

\begin{equation}
\mathcal{L}_{\text{CITA}} = \mathcal{L}_{\text{SFT}} + \lambda_1 \mathcal{L}_{\text{DPO}} + \lambda_2 \mathcal{L}_{\text{KL}}
\label{eq:unified_loss}
\end{equation}

where $\lambda_1, \lambda_2 > 0$ are balancing coefficients. We now detail each component:

\paragraph{SFT Loss (Response Quality).}
Standard cross-entropy loss ensuring high-quality response generation:

\begin{equation}
\mathcal{L}_{\text{SFT}} = -\frac{1}{N} \sum_{(x, y^+)} \sum_{t=1}^{T} \log P_\theta(y_t^+ \mid x, y_{<t}^+)
\label{eq:sft_loss}
\end{equation}

where $x$ is the instruction-augmented prompt, $y^+$ is the preferred response, and $T$ is the sequence length.

\paragraph{DPO Loss (Preference Alignment).}
Contrastive preference optimization without a separate reward model~\cite{rafailov2023direct,ziegler2019fine}:

\begin{equation}
\mathcal{L}_{\text{DPO}} = -\frac{1}{N} \sum \log \sigma \left( \beta \cdot \Delta_\theta \right)
\label{eq:dpo_loss}
\end{equation}

where the log-probability difference is:
\begin{equation}
\Delta_\theta = \log P_\theta(y^+ \mid x) - \log P_\theta(y^- \mid x)
\label{eq:delta}
\end{equation}

and $\beta > 0$ is the temperature parameter controlling preference sharpness.

\paragraph{KL Loss (Stability --- Mandatory).}
The critical component enabling behavioral switching, inspired by KL-regularized RL~\cite{jaques2019way,korbak2022rl} and continual learning~\cite{kirkpatrick2017overcoming,li2017learning}:

\begin{equation}
\mathcal{L}_{\text{KL}} = \frac{1}{N} \sum_{(x,y)} \text{KL}\left[ P_\theta(\cdot \mid x) \,\|\, P_{\text{ref}}(\cdot \mid x) \right]
\label{eq:kl_loss}
\end{equation}

where $P_{\text{ref}}$ is the reference model (pre-CITA checkpoint).

\vspace{-2mm}
\subsection{Why Mandatory KL?}
\label{sec:mandatory_kl}

\begin{table}[H]
\centering
\small
\begin{tabular}{@{}lcc@{}}
\toprule
\textbf{Scenario} & \textbf{Without KL} & \textbf{With KL} \\
\midrule
Mode Collapse & High risk & Prevented \\
Instruction Switching & Unstable & Stable \\
Behavioral Consistency & Degrades & Maintained \\
Preference Margin & Unbounded & Controlled \\
\bottomrule
\end{tabular}
\caption{Effect of mandatory KL regularization on training stability.}
\label{tab:kl_effects}
\vspace{-4mm}
\end{table}

\textbf{Key Insight.} In standard DPO, KL regularization is implicit and optional. For CITA, we make it \textit{explicit and mandatory} because:

\begin{enumerate}[leftmargin=*,itemsep=2pt]
    \item \textbf{Prevents Mode Collapse}: Without KL, the model can overfit to specific instruction patterns, losing generalization---a form of catastrophic forgetting~\cite{french1999catastrophic,kirkpatrick2017overcoming}.

    \item \textbf{Enables Switching}: KL keeps the model ``close enough'' to the base policy to respond appropriately to different instructions rather than collapsing to a single behavioral mode, similar to knowledge distillation~\cite{hinton2015distilling}.

    \item \textbf{Bounds Preference Margins}: Unconstrained optimization can lead to extreme log-probability differences, destabilizing training---a known issue in PPO~\cite{schulman2017proximal}, DPO variants~\cite{azar2023general,zhao2023slic}, and gradient-based continual learning~\cite{lopez2017gradient}.
\end{enumerate}

\vspace{-2mm}
\subsection{Gradient Analysis}
\label{sec:gradient}

The gradient of the contrastive term with respect to model parameters is:

\begin{equation}
\nabla_\theta \mathcal{L}_{\text{DPO}} = -\beta \left[ (1 - P^+) \nabla_\theta s^+ - P^+ \nabla_\theta s^- \right]
\label{eq:gradient}
\end{equation}

where $P^+ = \sigma(\beta \cdot \Delta_\theta)$ is the preference probability and $s^{\pm} = \log P_\theta(y^{\pm} \mid x)$.

\textbf{Interpretation.} When $P^+ \approx 1$ (model correctly prefers $y^+$), gradients become small, preventing over-optimization. The KL term adds a stabilizing force that prevents $P^+$ from saturating too quickly.
